\ifdefined\gkver\else\def\gkver{student}\fi
\documentclass[\gkver]{gaokaozhenti}
\begin{document}

\chapter{2019年4月浙江}

\begin{enumerate}
\item
%% number: 1
%% paperName: 2019年4月浙江选考第 1 题 3 分
%% typeId: 1
%% type: 单选题
%% chapter: 运动和力
%% point: 力学单位制
%% method: 无
%% score: 3
%% degree: 950
%% duplicateId: 0
%% body:
下列物理量属于基本量且单位属于国际单位制中基本单位的是\xzanswer{B}
\fourchoices[answer=B]
{功/焦耳}
{质量/千克}
{电荷量/库仑}
{力/牛顿}
%% answer: B
%% memo:
\memoanswer{
属于国际单位制中基本单位的是 kg（千克）、m（米）、s（秒）、A（安培）、mol（摩尔）、cd（坎德拉）、K（开尔文），对应的基本物理量是质量、长度、时间、电流、物质的量、发光强度、热力学温度，牛顿、焦耳、库仑是导出单位，故 B 正确 ACD 错误。
}

\item
%% number: 2
%% paperName: 2019年4月浙江选考第 2 题 3 分
%% typeId: 1
%% type: 单选题
%% chapter: 电路
%% point: 电容
%% method: 无
%% score: 3
%% degree: 950
%% duplicateId: 0
%% body:
下列器件中是电容器的是\xzanswer{B}
\fourchoices[answer=B, ispicture=true, h=2.0cm]
{figs/02a.png}
{figs/02b.png}
{figs/02c.png}
{figs/02d.png}
%% answer: B
%% memo:
\memoanswer{
A 图为滑动变阻器，B 图为电容器，C 图为电阻箱，D 图为定值电阻，故 B 正确 ACD 错误。
}

\item
%% number: 3
%% paperName: 2019年4月浙江选考第 3 题 3 分
%% typeId: 1
%% type: 单选题
%% chapter: 其他
%% point: 物理思想
%% method: 无
%% score: 3
%% degree: 750
%% duplicateId: 0
%% body:
下列式子属于比值定义物理量的是\xzanswer{C}
\fourchoices[answer=C]
{$t = \frac{{\Delta x}}{v}$}
{$a = \frac{F}{m}$}
{$C = \frac{Q}{U}$}
{$I = \frac{U}{R}$}
%% answer: C
%% memo:
\memoanswer{
电容 $C=\frac{Q}{U}$ 是比值定义式，加速度 $a=\frac{F}{m}$ 是决定式，电流 $I=\frac{U}{R}$ 是决定式，时间 $t=\frac{\Delta x}{v}$ 是计算式，既不是比值定义式，也不是决定式，故 C 正确 ABD 错误。
}

\item
%% number: 4
%% paperName: 2019年4月浙江选考第 4 题 3 分
%% typeId: 1
%% type: 单选题
%% chapter: 其他
%% point: 物理学史
%% method: 无
%% score: 3
%% degree: 850
%% duplicateId: 0
%% body:
下列陈述与事实相符的是\xzanswer{D}
\fourchoices[answer=D]
{牛顿测定了引力常量}
{法拉第发现了电流周围存在磁场}
{安培发现了静电荷间的相互作用规律}
{伽利略指出了力不是维持物体运动的原因}
%% answer: D
%% memo:
\memoanswer{
牛顿发现万有引力定律，卡文迪许最早测定了引力常量，故 A 错误；奥斯特最早发现了电流周围存在磁场，故 B 错误；库仑最早发现了静止的点电荷之间的相互作用规律，故 C 错误；伽利略通过理想斜面实验，最早指出了力不是维持物体运动的原因，故 D 正确。
}

\item
%% number: 5
%% paperName: 2019年4月浙江选考第 5 题 3 分
%% typeId: 1
%% type: 单选题
%% chapter: 磁场
%% point: 安培力
%% method: 无
%% score: 3
%% degree: 850
%% duplicateId: 0
%% body:
在磁场中的同一位置放置一条直导线，导线的方向与磁场方向垂直，则下列描述导线受到的安培力 $F$ 的大小与通过导线的电流 $I$ 的关系图象正确的是\xzanswer{A}
\fourchoices[answer=A, ispicture=true, h=2.4cm]
{figs/05a.png}
{figs/05b.png}
{figs/05c.png}
{figs/05d.png}
%% answer: A
%% memo:
\memoanswer{
在磁场的同一位置，磁感应强度 $B$ 保持恒定不变。当导线方向与磁场方向垂直时，安培力大小 $F=BIL$，所以安培力正比于电流 $I$，故 A 正确 BCD 错误。
}

\item
%% number: 6
%% paperName: 2019年4月浙江选考第 6 题 3 分
%% typeId: 1
%% type: 单选题
%% chapter: 运动和力
%% point: 牛顿第三定律
%% method: 无
%% score: 3
%% degree: 850
%% duplicateId: 0
%% body:
如图所示，小明撑杆使船离岸，则下列说法正确的是\xzanswer{A}
\onepicture[h=2.0cm, align=right]{figs/06.jpg}
\fourchoices[answer=A]
{小明与船之间存在摩擦力}
{杆的弯曲是由于受到杆对小明的力}
{杆对岸的力大于岸对杆的力}
{小明对杆的力和岸对杆的力是一对相互作用力}
%% answer: A
%% memo:
\memoanswer{
若小明与船之间没有摩擦力，小明将与船之间产生相对滑动，船也将不会远离河岸，所以小明和船之间一定存在静摩擦力，故 A 正确；杆的弯曲是由于杆受到了小明的作用力，故 B 错误；杆对岸的力与岸对杆的力是一对相互作用力，由牛顿第三定律知，两者大小相等，故 C 错误；小明对杆的力与岸对杆的力，作用在同一个物体（杆）上，不是一对相互作用力，故 D 错误。
}

\item
%% number: 7
%% paperName: 2019年4月浙江选考第 7 题 3 分
%% typeId: 1
%% type: 单选题
%% chapter: 曲线运动
%% point: 人造地球卫星
%% method: 无
%% score: 3
%% degree: 650
%% duplicateId: 0
%% body:
某颗北斗导航卫星属于地球静止轨道卫星（即卫星相对于地面静止）。则此卫星的\xzanswer{C}
\onepicture[h=2.6cm, align=right]{figs/07.jpg}
\fourchoices[answer=C]
{线速度大于第一宇宙速度}
{周期小于同步卫星的周期}
{角速度大于月球绕地球运行的角速度}
{向心加速度大于地面的重力加速度}
%% answer: C
%% memo:
\memoanswer{
对于地球卫星，由万有引力提供向心力有 $G\frac{Mm}{r^{2}}=m\frac{v^{2}}{r}$，得线速度为 $v=\sqrt{\frac{GM}{r}}$，可知卫星的运行半径越大，线速度越小，因该卫星的轨道半径大于地球半径，故其线速度小于近地卫星的线速度，即小于第一宇宙速度，故 A 错误；因该卫星相对于地面静止，周期应等于地球自转周期，即等于同步卫星的周期，故 B 错误；因该卫星周期为 $24\Uh=1$ 天，小于月球绕地球转动周期约 29 天，由角速度公式 $\omega=\frac{2\pi}{T}$，可知该卫星角速度较大，故 C 正确；根据 $G\frac{Mm}{r^2}=ma$，得向心加速度 $a=\frac{GM}{r^2}$，因该卫星的运行半径大于地球半径，故其向心加速度小于近地卫星的向心加速度，即小于地面的重力加速度 $g=\frac{GM}{R^2}$，故 D 错误。
}

\item
%% number: 8
%% paperName: 2019年4月浙江选考第 8 题 3 分
%% typeId: 1
%% type: 单选题
%% chapter: 电路
%% point: 非纯电阻电路
%% method: 无
%% score: 3
%% degree: 750
%% duplicateId: 0
%% body:
电动机与小电珠串联接入电路，电动机正常工作时，小电珠的电阻为 $R_{1}$，两端电压为 $U_{1}$，流过的电流为 $I_{1}$；电动机的内电阻为 $R_{2}$，两端电压为 $U_{2}$，流过的电流为 $I_{2}$。则\xzanswer{D}
\fourchoices[answer=D]
{$I_{1} < I_{2}$}
{$\frac{{{U_1}}}{{{U_2}}} > \frac{{{R_1}}}{{{R_2}}}$}
{$\frac{{{U_1}}}{{{U_2}}} = \frac{{{R_1}}}{{{R_2}}}$}
{$\frac{{{U_1}}}{{{U_2}}} < \frac{{{R_1}}}{{{R_2}}}$}
%% answer: D
%% memo:
\memoanswer{
因为两者串联，可知 $I_{1}=I_{2}$，选项 A 错误；因电动机为非纯电阻电器，有 $U_{2}>I_{2}R_{2}$，而小电珠为纯电阻，有 $U_{1}=I_{1}R_{1}$，则 $\frac{U_{1}}{U_{2}}<\frac{R_{1}}{R_{2}}$，选项 B、C 错误；对于电动机，电功率等于输出的机械功率与热功率之和，即 $U_{2}I_{2}=P_{\text{机械}}+I_{2}^{2}R_{2}$，则有 $U_{2}I_{2}>I_{2}^{2}R_{2}$，得 $U_{2}>I_{2}R_{2}$，综上所述 $\frac{U_{1}}{U_{2}}<\frac{R_{1}}{R_{2}}$，故 D 正确。
}

\item
%% number: 9
%% paperName: 2019年4月浙江选考第 9 题 3 分
%% typeId: 1
%% type: 单选题
%% chapter: 直线运动
%% point: 运动图象
%% method: 无
%% score: 3
%% degree: 750
%% duplicateId: 0
%% body:
甲、乙两物体零时刻开始从同一地点向同一方向做直线运动，位移-时间图象如图所示，则在 $0\sim t_{1}$ 时间内\xzanswer{B}
\onepicture[h=2.6cm, align=right]{\includesvg{figs/09}}
\fourchoices[answer=B]
{甲的速度总比乙大}
{甲、乙位移相同}
{甲经过的路程比乙小}
{甲、乙均做加速运动}
%% answer: B
%% memo:
\memoanswer{
在 $x-t$ 图象中，斜率代表物体速度大小，由题图可知在 $t_{1}$ 时刻，乙的速度大于甲的速度，故 A 错误；甲、乙初始位置相同，末位置相同，二者位移相同，故 B 正确；甲、乙从同一地点向同一方向做直线运动，始末位置相同，故路程相等，故 C 错误；甲的斜率不变，甲做匀速直线运动，乙的斜率变大，乙做加速直线运动，故 D 错误。
}

\item
%% number: 10
%% paperName: 2019年4月浙江选考第 10 题 3 分
%% typeId: 1
%% type: 单选题
%% chapter: 电场
%% point: 带电粒子的直线运动
%% method: 无
%% score: 3
%% degree: 550
%% duplicateId: 0
%% body:
当今医学上对某些肿瘤采用质子疗法进行治疗，该疗法用一定能量的质子束照射肿瘤杀死癌细胞。现用一直线加速器来加速质子，使其从静止开始被加速到 $1.0\times10^{7}\Ums$。已知加速电场的场强为 $1.3\times10^{5}\UNC$，质子的质量为 $1.67\times10^{-27}\Ukg$，电荷量为 $1.6\times10^{-19}\UC$，则下列说法正确的是\xzanswer{D}
\onepicture[h=2.4cm, align=right]{figs/10.jpg}
\fourchoices[answer=D]
{加速过程中质子电势能增加}
{质子所受到的电场力约为 $2\times10^{-15}\UN$}
{质子加速需要的时间约为 $8\times10^{-6}\Us$}
{加速器加速的直线长度约为 $4\Um$}
%% answer: D
%% memo:
\memoanswer{
质子在电场中加速过程，电场力对其做正功，质子电势能减少，故 A 错误；质子受到的电场力为 $F=qE=1.3\times10^{5}\times1.6\times10^{-19}\UN=2.08\times10^{-14}\UN$，故 B 错误；由牛顿第二定律得，质子加速度 $a=\frac{F}{m}=\frac{2.08\times10^{-14}}{1.67\times10^{-27}}\Umsq=1.25\times10^{13}\Umsq$，由匀变速运动公式 $v=v_{0}+at$ 得，加速时间 $t=8\times10^{-7}\Us$，故 C 错误；由 $v^{2}-v_{0}^{2}=2ax$，得加速直线长度 $x=4\Um$，故 D 正确。
}

\item
%% number: 11
%% paperName: 2019年4月浙江选考第 11 题 3 分
%% typeId: 1
%% type: 单选题
%% chapter: 力物体的平衡
%% point: 物体系平衡
%% method: 整体与隔离
%% score: 3
%% degree: 450
%% duplicateId: 0
%% body:
如图所示，一根粗糙的水平横杆上套有 A、B 两个轻环，系在两环上的等长细绳拴住的书本处于静止状态，现将两环距离变小后书本仍处于静止状态，则\xzanswer{B}
\onepicture[h=2.6cm, align=right]{figs/11.png}
\fourchoices[answer=B]
{杆对 A 环的支持力变大}
{B 环对杆的摩擦力变小}
{杆对 A 环的力不变}
{与 B 环相连的细绳对书本的拉力变大}
%% answer: B
%% memo:
\memoanswer{
对两环和书本的整体受力分析，竖直方向：$2N=mg$，可知将两环距离变小后杆对 A 环的支持力不变，选项 A 错误；对圆环 B 受力分析可知，$f=T\cos\theta$；对书本：$2T\sin\theta=mg$，解得 $f=\frac{mg}{2\tan\theta}$（其中的 $\theta$ 是绳与杆之间的夹角），则当两环距离变小后，$\theta$ 变大，则 $f$ 减小，与 B 环相连的细绳对书本的拉力 $T$ 变小，选项 B 正确，D 错误；同理，杆对 A 环的摩擦力减小，杆对 A 环的支持力不变，则杆对 A 环的力减小，选项 C 错误；故选 B。
}
\memoanswer[另解]{
如图甲所示，设书本的质量为 $m$，将 A、B 两轻环、细绳与书本视为一整体，整体受到竖直向下的重力、杆对 A、B 环的竖直向上的支持力，两个支持力大小之和等于书本重力 $mg$，每个支持力大小为 $F_{N}=\frac{1}{2}mg$，与 A、B 环间距无关，故 A 错误；对 B 环受力分析如图乙所示，B 环受到细绳拉力 $F_{T}$、杆对 B 环向上的支持力 $F_{N}$ 及向右的摩擦力 $F_{f}$，由力的合成或分解及平衡条件，得 $F_{f}=F_{N}\tan\theta=\frac{1}{2}mg\tan\theta$，因两环间距变小，$\tan\theta$ 随着 $\theta$ 的变小而变小，故杆对 B 环的摩擦力 $F_{f}$ 变小，由牛顿第三定律知，B 环对杆的摩擦力变小，故 B 正确；对 A 环受力分析与 B 环类似，杆对 A 环的力是支持力 $F_{N}$ 和摩擦力 $F_{f}$ 的合力，该合力与细绳拉力 $F_{T}$ 等大、反向，由力的合成及平衡条件得 $F_{T}=\frac{F_{N}}{\cos\theta}=\frac{mg}{2\cos\theta}$，当 $\theta$ 变小时，$\cos\theta$ 变大，则 $F_{T}$ 变小，即杆对 A 环的力变小，故 C 错误；如图丙所示，对书本进行受力分析，由力的合成与平衡条件，得 $2F_{T}\cos\theta=mg$，随两环间距变小，$\theta$ 变小，$\cos\theta$ 变大，故细绳拉力 $F_{T}$ 变小，故 D 错误。

\threepicture[h=3.0cm, showlabel=false]{figs/11_an_a.png}{figs/11_an_b.png}{figs/11_an_c.png}
}

\item
%% number: 12
%% paperName: 2019年4月浙江选考第 12 题 3 分
%% typeId: 1
%% type: 单选题
%% chapter: 运动和力
%% point: 超重失重
%% method: 无
%% score: 3
%% degree: 450
%% duplicateId: 0
%% body:
如图所示，A、B、C 为三个实心小球，A 为铁球，B、C 为木球。A、B 两球分别连在两根弹簧上，C 球连接在细线一端，弹簧和细线的下端固定在装水的杯子底部，该水杯置于用绳子悬挂的静止吊篮内。若将挂吊篮的绳子剪断，则剪断的瞬间相对于杯底（不计空气阻力，$\rho_{\text{木}} < \rho_{\text{水}} < \rho_{\text{铁}}$）\xzanswer{D}
\onepicture[h=3.0cm, align=right]{figs/12.png}
\fourchoices[answer=D]
{A 球将向上运动，B、C 球将向下运动}
{A、B 球将向上运动，C 球不动}
{A 球将向下运动，B 球将向上运动，C 球不动}
{A 球将向上运动，B 球将向下运动，C 球不动}
%% answer: D
%% memo:
\memoanswer{
剪断绳子前 A、B、C 三个实心小球在水杯中均为静止状态，由牛顿第二定律得，对 A 球 $F_{A\text{弹}}=m_{A}g-F_{A\text{浮}}$，对 B 球 $F_{B\text{弹}}=F_{B\text{浮}}-m_{B}g$，C 球 $F_{C\text{浮}}=m_{C}g+F_{T}$，将挂吊篮的绳子剪断瞬间，装有水的杯子仅受到重力作用，水处于完全失重状态，即可以认为水和球之间没有相互作用力，即原有的浮力突然变为 0，但弹簧弹力不发生突变。以杯底作为参考系，此时 A 球受到向上的弹簧弹力大于 A 球的重力，所以 A 球将向上运动；B 球受到向下的弹簧弹力与重力作用，此时 B 球将向下运动；C 不受浮力，与杯中的水一起自由落体运动，C 球相对杯底不动，故 D 正确 ABC 错误。
}

\item
%% number: 13
%% paperName: 2019年4月浙江选考第 13 题 3 分
%% typeId: 1
%% type: 单选题
%% chapter: 电场
%% point: 带电粒子的直线运动
%% method: 无
%% score: 3
%% degree: 350
%% duplicateId: 0
%% body:
用长为 $1.4\Um$ 的轻质柔软绝缘细线，拴一质量为 $1.0\times10^{-2}\Ukg$、电荷量为 $2.0\times10^{-8}\UC$ 的小球，细线的上端固定于 O 点。现加一水平向右的匀强电场，平衡时细线与铅垂线成 $37^\circ$，如图所示。现向左拉小球使细线水平且拉直，静止释放，则（$\sin37^\circ = 0.6$）\xzanswer{C}
\onepicture[h=2.6cm, align=right]{figs/13.png}
\fourchoices[answer=C]
{该匀强电场的场强为 $3.75\times10^{7}\UNC$}
{平衡时细线的拉力为 $0.17\UN$}
{经过 $0.5\Us$，小球的速度大小为 $6.25\Ums$}
{小球第一次通过 O 点正下方时，速度大小为 $7\Ums$}
%% answer: C
%% memo:
\memoanswer{
小球在匀强电场中处于平衡状态时，对小球进行受力分析，如图甲所示，由力的合成与平衡条件得，$qE=mg\tan37^\circ$，解得场强 $E=\frac{mg\tan37^\circ}{q}=3.75\times10^{6}\UNC$，故 A 错误；平衡时，细线拉力 $F_{T}=\frac{mg}{\cos37^\circ}=0.125\UN$，故 B 错误；向左拉小球使细线水平拉直，由静止释放后小球将沿与竖直线成 $37^\circ$ 角方向匀加速直线运动，如图乙所示，到达 B 点后绳子恰被拉直，此过程小球所受合力 $F_{\text{合}}=\frac{mg}{\cos37^\circ}$，加速度 $a=\frac{F_{\text{合}}}{m}=\frac{g}{\cos37^\circ}=12.5\Umsq$；在等腰三角形 OAB 中，AB 间距 $x_{AB}=2l\cos53^\circ=1.68\Um$，把 $t=0.5\Us$ 代入 $x=\frac{1}{2}at^{2}$ 得 $x=1.5625\Um<x_{AB}$，即经过 $0.5\Us$ 小球不会到达 B 点，故小球速度 $v=at=6.25\Ums$，故 C 正确；当小球到达 B 点后，其速度沿绳（半径）方向的分速度突变为 0，即细线拉紧过程有机械能损失，小球第一次到达 O 点正下方时的速度，一定小于机械能守恒时到达最低点的速度，机械能守恒时有 $mgl+qEl=\frac{1}{2}mv'^{2}-0$，解得 $v'=7\Ums$，故 D 错误。

\twopicture[h=3.0cm, showlabel=false]{figs/13_an_a.png}{figs/13_an_b.png}
}

\item
%% number: 14
%% paperName: 2019年4月浙江选考第 14 题 2 分
%% typeId: 2
%% type: 多选题
%% chapter: 光学
%% point: 光的干涉
%% method: 无
%% score: 2
%% degree: 550
%% duplicateId: 0
%% body:
波长为 $\lambda_{1}$ 和 $\lambda_{2}$ 的两束可见光入射到双缝，在光屏上观察到干涉条纹，其中波长为 $\lambda_{1}$ 的光的条纹间距大于波长为 $\lambda_{2}$ 的条纹间距。则（下列表述中，脚标“1”和“2”分别代表波长为 $\lambda_{1}$ 和 $\lambda_{2}$ 的光所对应的物理量）\xzanswer{BD}
\fourchoices[answer=BD]
{这两束光的光子的动量 $p_{1} > p_{2}$}
{这两束光从玻璃射向真空时，其临界角 $C_{1} > C_{2}$}
{这两束光都能使某种金属发生光电效应，则遏止电压 $U_{1} > U_{2}$}
{这两束光由氢原子从不同激发态跃迁到 $n =$ 2 能级时产生，则相应激发态的电离能 $\Delta E_{1} > \Delta E_{2}$}
%% answer: BD
%% memo:
\memoanswer{
由双缝干涉间距公式 $\Delta x=\frac{l}{d}\lambda$ 可得 $\lambda_{1}>\lambda_{2}$，由光子动量公式 $p=\frac{h}{\lambda}$，可得 $p_{1}<p_{2}$，故 A 错误；由波速公式 $c=\lambda\nu$，可得光子频率 $\nu_{1}<\nu_{2}$，在同一玻璃中的折射率 $n_{1}<n_{2}$，由全反射临界角公式 $\sin C=\frac{1}{n}$，可得 $C_{1}>C_{2}$，故 B 正确；由光电效应方程 $h\nu=W+E_{k}$，结合 $E_{k}=eU_{c}$，可得遏止电压 $U_{c}=\frac{h}{e}\nu-\frac{W}{e}$，则有 $U_{c1}<U_{c2}$，故 C 错误；由光子能量公式 $E=h\nu=h\frac{c}{\lambda}$，可得波长为 $\lambda_{1}$ 的光所处的激发态能级 $E_{1}$ 更低，则从激发态跃迁到无穷远处需要的电离能 $\Delta E_{1}>\Delta E_{2}$，故 D 正确。
}

\item
%% number: 15
%% paperName: 2019年4月浙江选考第 15 题 2 分
%% typeId: 2
%% type: 多选题
%% chapter: 原子和原子核
%% point: 结合能
%% method: 无
%% score: 2
%% degree: 550
%% duplicateId: 0
%% body:
静止在匀强磁场中的原子核 X 发生 $\alpha$ 衰变后变成新原子核 Y。已知核 X 的质量数为 $A$，电荷数为 $Z$，核 X、核 Y 和 $\alpha$ 粒子的质量分别为 $m_{X}$、$m_{Y}$ 和 $m_{\alpha}$，$\alpha$ 粒子在磁场中运动的半径为 $R$。则\xzanswer{AC}
\fourchoices[answer=AC]
{衰变方程可表示为 $\ce{^{A}_{Z}X -> ^{A-4}_{Z-2}Y + ^{4}_{2}He}$}
{核 Y 的结合能为（$m_{X} - m_{Y} - m_{\alpha}$）$c^{2}$}
{核 Y 在磁场中运动的半径为 $\frac{{2R}}{{Z - 2}}$}
{核Y的动能为 $E_{kY} = \frac{{{m_Y}({m_X} - {m_Y} - {m_\alpha }){c^2}}}{{{m_Y} + {m_\alpha }}}$}
%% answer: AC
%% memo:
\memoanswer{
核反应前后质量数守恒和电荷数守恒，衰变方程可表示为 $\ce{^{A}_{Z}X -> ^{A-4}_{Z-2}Y + ^{4}_{2}He}$，故 A 正确；衰变过程因质量亏损而释放的核能为 $E=\Delta mc^{2}=(m_{X}-m_{Y}-m_{\alpha})c^{2}$，此式不是核 Y 的结合能，故 B 错误；核反应前后，系统动量守恒，即 $m_{Y}v_{Y}=m_{\alpha}v_{\alpha}$，新核在磁场中圆周运动的半径 $r=\frac{mv}{qB}$，所以半径正比于 $\frac{1}{q}$，即有 $\frac{r_{Y}}{R}=\frac{2}{Z-2}$，得 $r_{Y}=\frac{2R}{Z-2}$，故 C 正确；核反应过程，系统能量守恒，即释放核能等于 Y 核与 $\alpha$ 粒子的动能之和，则有 $(m_{X}-m_{Y}-m_{\alpha})c^{2}=\frac{1}{2}m_{Y}v_{Y}^{2}+\frac{1}{2}m_{\alpha}v_{\alpha}^{2}$，动能之比 $\frac{\frac{1}{2}m_{Y}v_{Y}^{2}}{\frac{1}{2}m_{\alpha}v_{\alpha}^{2}}=\frac{m_{Y}v_{Y}}{m_{\alpha}v_{\alpha}}\cdot\frac{v_{Y}}{v_{\alpha}}=\frac{v_{Y}}{v_{\alpha}}=\frac{m_{\alpha}}{m_{Y}}$，得 Y 核动能 $\frac{1}{2}m_{Y}v_{Y}^{2}=\frac{m_{\alpha}}{m_{Y}+m_{\alpha}}(m_{X}-m_{Y}-m_{\alpha})c^{2}$，故 D 错误。
}

\item
%% number: 16
%% paperName: 2019年4月浙江选考第 16 题 2 分
%% typeId: 2
%% type: 多选题
%% chapter: 机械波
%% point: 波与振动图象
%% method: 无
%% score: 2
%% degree: 450
%% duplicateId: 0
%% body:
图 1 为一列简谐横波在 $t$ = 0 时刻的波形图，P、Q 为介质中的两个质点，图 2 为质点 P 的振动图象，则\xzanswer{CD}
\twopicture[numA=1, numB=2, labelA=20194月浙江16a, labelB=20194月浙江16b, h=2.2cm, align=right]{figs/16a.png}{figs/16b.png}
\fourchoices[answer=CD]
{$t=0.2\Us$ 时，质点 Q 沿 $y$ 轴负方向运动}
{$0\sim0.3\Us$ 内，质点 Q 运动的路程为 $0.3\Um$}
{$t=0.5\Us$ 时，质点 Q 的加速度小于质点 P 的加速度}
{$t=0.7\Us$ 时，质点 Q 距平衡位置的距离小于质点 P 距平衡位置的距离}
%% answer: CD
%% memo:
\memoanswer{
由题图可知，波长 $\lambda=2\Um$，振幅 $A=0.1\Um$，周期 $T=0.4\Us$，由题图 2 可知质点 P 在 $t=0$ 时开始向 $y$ 轴正方向运动，即波在向 $x$ 轴负方向传播，质点 Q 在 $t=0$ 时向 $y$ 轴负方向运动，$t=0.2\Us=\frac{T}{2}$ 时，质点 Q 向 $y$ 轴正方向运动，故 A 错误；$0\sim0.3\Us$ 内，质点 Q 完成 $\frac{3}{4}T$ 振动，路程不是 3 倍振幅，即不是 $0.3\Um$，故 B 错误；质点 Q、P 在 $t=0.5\Us$ 和 $t=0.1\Us$ 时的位移相同，$t=0.1\Us$ 时质点 Q 未到最大位移处，$t=0.1\Us$ 时质点 P 位于波峰位置，其加速度最大，故 C 正确；质点 Q、P 在 $t=0.7\Us$ 与 $t=0.3\Us$ 时的位移相同，$t=0.3\Us$ 时 P 点位于波谷位置，距平衡位置的距离最大，$t=0.3\Us$ 时质点 Q 未到达最大位移处，故 D 正确。
}

\item
%% number: 17
%% paperName: 2019年4月浙江选考第 17 题 5 分
%% typeId: 4
%% type: 实验题
%% chapter: 曲线运动
%% point: 研究平抛运动实验
%% method: 无
%% score: 5
%% degree: 550
%% duplicateId: 0
%% body:
采用如图 1 所示的实验装置做“研究平抛运动”的实验。
\begin{enumerate}
	\item
	实验时需要下列哪个器材\tkanswer[1.5cm]{B}。

	A．弹簧秤\qquad B．重锤线\qquad C．打点计时器

	\item
	做实验时，让小球多次沿同一轨道运动，通过描点法画出小球平抛运动的轨迹。下列的一些操作要求，正确的是\tkanswer[1.5cm]{ACD}（多选）。

	A．每次必须由同一位置静止释放小球

	B．每次必须严格地等距离下降记录小球位置

	C．小球运动时不应与木板上的白纸相接触

	D．记录的点应适当多一些

	\item
	若用频闪摄影方法来验证小球在平抛过程中水平方向是匀速运动，记录下如图 2 所示的频闪照片。在测得 $x_{1}$、$x_{2}$、$x_{3}$、$x_{4}$ 后，需要验证的关系是\tkanswer[3.5cm]{$x_{4}-x_{3}=x_{3}-x_{2}=x_{2}-x_{1}=x_{1}$}。已知频闪周期为 $T$，用下列计算式求得的水平速度，误差较小的是\tkanswer[1cm]{D}。

	A．$\frac{x_{1}}{T}$\qquad B．$\frac{x_{2}}{2T}$\qquad C．$\frac{x_{3}}{3T}$\qquad D．$\frac{x_{4}}{4T}$
\end{enumerate}
\twopicture[numA=1, numB=2, labelA=20194月浙江17a, labelB=20194月浙江17b, h=3.0cm, align=right]{figs/17a.png}{\includesvg{figs/17b}}
%% answer: 
\jdanswer{
\begin{enumerate}
	\item
	B

	\item
	ACD

	\item
	$x_{4}-x_{3}=x_{3}-x_{2}=x_{2}-x_{1}=x_{1}$；D
\end{enumerate}
}
%% memo:
\memoanswer{
（1）该实验需要使用重垂线确定 $y$ 轴位置，故 B 正确；弹簧秤和打点计时器不需要使用，故 A、C 错误。

（2）每次必须由同一位置静止释放小球，保证每次轨迹都相同，A 正确；为了提高实验精度，使描绘的曲线更接近于平抛运动真实的轨迹，应尽量多地记录点，但不需要每次等距离下降记录小球的位置，故 B 错误、D 正确；小球运动过程如果和白纸相接触，就会有摩擦阻力的影响，小球轨迹不再是平抛运动的轨迹，所以小球不应与白纸相接触，故 C 正确。

（3）若小球在水平方向做匀速直线运动，则相邻的两张照片间的水平间距应相等，即 $x_{1}=x_{2}-x_{1}=x_{3}-x_{2}=x_{4}-x_{3}$；选用间距较远的两点间的水平位移 $x_{4}$ 进行水平初速度的计算 $v_{0}=\frac{x_{4}}{4T}$，误差会较小，故 D 正确。
}

\item
%% number: 18
%% paperName: 2019年4月浙江选考第 18 题 5 分
%% typeId: 4
%% type: 实验题
%% chapter: 电路
%% point: 伏安法测电阻
%% method: 无
%% score: 5
%% degree: 550
%% duplicateId: 0
%% body:
小明想测额定电压为 $2.5\UV$ 的小灯泡在不同电压下的电功率，设计了如图 1 所示的电路。
\begin{enumerate}
	\item
	在实验过程中，调节滑片 P，电压表和电流表均有示数但总是调不到零，其原因是\tkanswer[2.5cm]{1 点至 4 点}导线没有连接好（图中用数字标记的小圆点表示接线点，空格中请填写图中的数字，如“7 点至 8 点”）；

	\item
	正确连好电路，闭合开关，调节滑片 P，当电压表的示数达到额定电压时，电流表的指针如图 2 所示，则电流为\tkanswer[1.2cm]{0.30} $\UA$，此时小灯泡的功率为\tkanswer[1.2cm]{0.75} $\UW$；

	\item
	做完实验后小明发现在实验报告上漏写了电压为 $1.00\UV$ 时通过小灯泡的电流，但在草稿纸上记录了下列数据，你认为最有可能的是\tkanswer[1cm]{C}。

	（A）$0.08\UA$\qquad （B）$0.12\UA$\qquad （C）$0.20\UA$
\end{enumerate}
\twopicture[numA=1, numB=2, labelA=20194月浙江18a, labelB=20194月浙江18b, h=3.0cm, align=right]{figs/18a.png}{figs/18b.png}
%% answer: 
\jdanswer{
\begin{enumerate}
	\item
	1 点至 4 点

	\item
	0.30，0.75

	\item
	C
\end{enumerate}
}
%% memo:
\memoanswer{
（1）1 点到 4 点未连接好时，滑动变阻器处于限流接法，电压和电流均有示数且不能同时调节为 0。

（2）由题图 2 可知电流为 $0.30\UA$；小灯泡功率 $P=UI=2.5\times0.30\UW=0.75\UW$。

（3）当灯泡正常发光即电压为 $2.5\UV$ 时，小灯泡电阻为 $R=\frac{U}{I}=\frac{2.5}{0.3}\UO\approx8.3\UO$，由于灯丝电阻随温度升高而增大，故当其电压为 $1.0\UV$ 时，小灯泡电阻小于 $8.3\UO$，则有 $\frac{1}{I}<8.3$，电流 $I>\frac{1}{8.3}\UA\approx0.12\UA$，即 C 正确。
}

\item
%% number: 19
%% paperName: 2019年4月浙江选考第 19 题 9 分
%% typeId: 6
%% type: 计算题
%% chapter: 运动和力
%% point: 牛顿定律的应用
%% method: 无
%% score: 9
%% degree: 450
%% duplicateId: 0
%% body:
小明以初速度 $v_{0}=10\Ums$ 竖直向上抛出一个质量 $m=0.1\Ukg$ 的小皮球，最后在抛出点接住。假设小皮球在空气中所受阻力大小为重力的 0.1 倍。求小皮球
\begin{enumerate}
	\item
	上升的最大高度；
	\item
	从抛出到接住的过程中重力和空气阻力所做的功
	\item
	上升和下降的时间。
\end{enumerate}
%% answer: 
\jdanswer{
\begin{enumerate}
	\item
	$\frac{50}{11}\Um$

	\item
	$0$；$-\frac{10}{11}\UJ$

	\item
	$\frac{10}{11}\Us$，$\frac{10\sqrt{11}}{33}\Us$
\end{enumerate}
}
%% memo:
\memoanswer{
（1）小皮球在上升过程中，由牛顿第二定律得 $mg+F_{f}=ma_{1}$，解得 $a_{1}=11\Umsq$。小皮球上升的高度，由运动学公式得 $h=\frac{v_{0}^{2}}{2a_{1}}=\frac{50}{11}\Um$。

（2）小皮球从抛出到接住的过程中，重力做的功为 $W_{G}=0$；空气阻力做的功为 $W_{f}=-F_{f}\cdot2h=-\frac{10}{11}\UJ$。

（3）小皮球上升的时间为 $t_{1}=\frac{v_{0}}{a_{1}}=\frac{10}{11}\Us$；小皮球在下降过程中，由牛顿第二定律得 $mg-F_{f}=ma_{2}$，解得 $a_{2}=9\Umsq$，由运动学公式得 $h=\frac{1}{2}a_{2}t_{2}^{2}$，解得 $t_{2}=\frac{10\sqrt{11}}{33}\Us$。
}

\item
%% number: 20
%% paperName: 2019年4月浙江选考第 20 题 12 分
%% typeId: 6
%% type: 计算题
%% chapter: 曲线运动
%% point: 平抛运动
%% method: 无
%% score: 12
%% degree: 450
%% duplicateId: 0
%% body:
某砂场为提高运输效率，研究砂粒下滑的高度与砂粒在传送带上运动的关系，建立如图所示的物理模型。竖直平面内有一倾角 $\theta=37^\circ$ 的直轨道 AB，其下方右侧放置一水平传送带，直轨道末端 B 与传送带间距可近似为零，但允许砂粒通过。转轮半径 $R=0.4\Um$、转轴间距 $L=2\Um$ 的传送带以恒定的线速度逆时针转动，转轮最低点离地面的高度 $H=2.2\Um$。现将一小物块放在距离传送带高 $h$ 处静止释放，假设小物块从直轨道 B 端运动到达传送带上 C 点时，速度大小不变，方向变为水平向右。已知小物块与直轨道和传送带间的动摩擦因数均为 $\mu=0.5$。（$\sin37^\circ = 0.6$）
\begin{enumerate}
	\item
	若 $h=2.4\Um$，求小物块到达 B 端时速度的大小；
	\item
	若小物块落到传送带左侧地面，求 $h$ 需要满足的条件
	\item
	改变小物块释放的高度 $h$，小物块从传送带的 D 点水平向右抛出，求小物块落地点到 D 点的水平距离 $x$ 与 $h$ 的关系式及 $h$ 需要满足的条件。
\end{enumerate}
\onepicture[h=4.2cm, align=right]{figs/20.png}
%% answer: 
\jdanswer{
\begin{enumerate}
	\item
	$4\Ums$

	\item
	$h<3.0\Um$

	\item
	$h\geq3.6\Um$ 时，$x=\sqrt{4h-12}\Um$
\end{enumerate}
}
%% memo:
\memoanswer{
（1）物块由静止释放到 B 的过程中，由牛顿第二定律得 $mg\sin\theta-\mu mg\cos\theta=ma$，解得 $a=2\Umsq$，由运动学公式得 $v_{B}=2a\frac{h}{\sin\theta}=4\Ums$。

（2）设当小物块到 D 点的速度恰好为零时，此时 A 点距传送带高度为 $h_{1}$，由动能定理得 $0=mgh_{1}-\mu mg\cos\theta\frac{h_{1}}{\sin\theta}-\mu mgL$，解得 $h_{1}=3.0\Um$，当 $h<h_{1}=3.0\Um$ 时，小物块能够落到传送带左侧地面。

（3）设小物块从传送带右侧抛出时，D 点的速度为 $v$，由动能定理得 $\frac{1}{2}mv^{2}=mgh-\mu mg\cos\theta\frac{h}{\sin\theta}-\mu mgL$，由平抛的规律知 $H+2R=\frac{1}{2}gt^{2}$，$x=vt$，联立解得 $x=\sqrt{4h-12}$；为使能在 D 点水平抛出，重力不足以提供向心力，有 $mg\leq\frac{mv^{2}}{R}$，解得 $h\geq3.6\Um$。
}

\item
%% number: 21
%% paperName: 2019年4月浙江选考第 21 题 4 分
%% typeId: 4
%% type: 实验题
%% chapter: 电磁感应
%% point: 验证电磁感应实验
%% method: 无
%% score: 4
%% degree: 550
%% duplicateId: 0
%% body:
在“探究电磁感应的产生条件”实验中，实验连线后如图 1 所示，感应线圈组的内外线圈的绕线方向如图 2 粗线所示。
\begin{enumerate}
	\item
	接通电源，闭合开关，G 表指针会有大的偏转，几秒后 G 表指针停在中间不动。将滑动变阻器的触头迅速向右滑动时，G 表指针\tkanswer[1.2cm]{左偏}（“不动”、“右偏”、“左偏”、“不停振动”）；迅速抽出铁芯时，G 表指针\tkanswer[1.2cm]{右偏}（“不动”、“右偏”、“左偏”、“不停振动”）。
	\item
	断开开关和电源，将铁芯重新插入内线圈中，把直流输出改为交流输出，其他均不变。接通电源，闭合开关，G 表指针\tkanswer[1.5cm]{不停振动}（“不动”、“右偏”、“左偏”、“不停振动”）。
	\item
	仅用一根导线，如何判断 G 表内部线圈是否断了？\tkanswer[4cm]{短接 G 表前后各摇动 G 表一次，比较指针偏转，有明显变化，则线圈未断；反之则断了。}
\end{enumerate}
\twopicture[numA=1, numB=2, labelA=20194月浙江21a, labelB=20194月浙江21b, h=2.4cm, align=right]{figs/21a.png}{figs/21b.png}
%% answer: 
\jdanswer{
\begin{enumerate}
	\item
	左偏，右偏

	\item
	不停振动

	\item
	短接 G 表前后各摇动 G 表一次，比较指针偏转，有明显变化，则线圈未断；反之则断了。
\end{enumerate}
}
%% memo:
\memoanswer{
（1）将滑动变阻器的触头迅速向右滑动时，电阻减小，回路电流变大，根据线圈中导线的绕向可知磁通量向下增加，根据楞次定律可知，A 线圈中产生的感应电流使 G 表指针左偏；迅速抽出铁芯时，磁通量减小，产生的感应电流方向与上述方向相反，则 G 表指针右偏。

（2）断开开关和电源，将铁芯重新插入内线圈中，把直流输出改为交流输出，其他均不变。接通电源，闭合开关，由于穿过线圈的磁通量大小方向都不断变化，在线圈 A 中产生的感应电流大小方向不断变化，则 G 表指针不停振动。

（3）根据电磁阻尼原理，短接 G 表，前后各摇动 G 表一次，比较指针偏转，有明显变化，则线圈未断；反之则断了。
}

\item
%% number: 22
%% paperName: 2019年4月浙江选考第 22 题 10 分
%% typeId: 6
%% type: 计算题
%% chapter: 电磁感应
%% point: 电磁感应中的变加速
%% method: 无
%% score: 10
%% degree: 350
%% duplicateId: 0
%% body:
如图所示，倾角 $\theta=37^\circ$、间距 $l=0.1\Um$ 的足够长金属导轨底端接有阻值 $R=0.1\UO$ 的电阻，质量 $m=0.1\Ukg$ 的金属棒 ab 垂直导轨放置，与导轨间的动摩擦因数 $\mu=0.45$。建立原点位于底端、方向沿导轨向上的坐标轴 $x$。在 $0.2\Um\leq x\leq0.8\Um$ 区间有垂直导轨平面向上的匀强磁场。从 $t=0$ 时刻起，棒 ab 在沿 $x$ 轴正方向的外力 $F$ 作用下从 $x=0$ 处由静止开始沿斜面向上运动，其速度与位移 $x$ 满足 $v=kx$（可导出 $a=kv$），$k=5\Usn$。当棒 ab 运动至 $x_{1}=0.2\Um$ 处时，电阻 $R$ 消耗的电功率 $P=0.12\UW$，运动至 $x_{2}=0.8\Um$ 处时撤去外力 $F$，此后棒 ab 将继续运动，最终返回至 $x=0$ 处。棒 ab 始终保持与导轨垂直，不计其它电阻，求：（提示：可以用 $F-x$ 图象下的“面积”代表力 $F$ 做的功）
\begin{enumerate}
	\item
	磁感应强度 $B$ 的大小
	\item
	外力 $F$ 随位移 $x$ 变化的关系式；
	\item
	在棒 ab 整个运动过程中，电阻 $R$ 产生的焦耳热 $Q$。
\end{enumerate}
\onepicture[h=3.4cm, align=right]{figs/22.png}
%% answer: 
\jdanswer{
\begin{enumerate}
	\item
	$\frac{\sqrt{30}}{5}\UT$

	\item
	无磁场区间：$F=0.96+2.5x$；有磁场区间：$F=0.96+3.1x$

	\item
	$0.324\UJ$
\end{enumerate}
}
%% memo:
\memoanswer{
（1）由功率的表达式知 $P=\frac{(Blv)^{2}}{R}$，解得 $B=\sqrt{\frac{PR}{(lv)^{2}}}=\frac{\sqrt{30}}{5}\UT$。

（2）①无磁场区间 $0\leq x<0.2\Um$，由题意知 $a=5v=25x$，由受力分析和牛顿第二定律可知外力与位移 $x$ 的关系为 $F=25xm+\mu mg\cos\theta+mg\sin\theta=0.96+2.5x$。②有磁场区间 $0.2\Um\leq x\leq0.8\Um$，金属棒受到的安培力为 $F_{A}=\frac{(Bl)^{2}v}{R}=0.6x$，由受力分析和牛顿第二定律可知外力与位移 $x$ 的关系为 $F=0.96+2.5x+0.6x=0.96+3.1x$。

（3）棒 ab 上升过程中克服安培力做功（梯形面积）为 $W_{A1}=\frac{0.6}{2}(x_{1}+x_{2})(x_{2}-x_{1})=0.18\UJ$；撤去外力后，棒 ab 上升的最大距离为 $s$，再次进入磁场时的速度为 $v'$，由动能定理得 $\frac{1}{2}mv^{2}=(mg\sin\theta+\mu mg\cos\theta)s$，$\frac{1}{2}mv'^{2}=(mg\sin\theta-\mu mg\cos\theta)s$，联立解得 $v'=2\Ums$；由于 $mg\sin\theta-\mu mg\cos\theta-\frac{(Bl)^{2}v'}{R}=0$，即金属棒受力平衡，故棒 ab 再次进入磁场后做匀速运动。下降过程中克服安培力做功为 $W_{A2}=\frac{(Bl)^{2}v'}{R}(x_{2}-x_{1})=0.144\UJ$，故电阻 R 产生的总焦耳热为 $Q=W_{A1}+W_{A2}=0.324\UJ$。
}

\item
%% number: 23
%% paperName: 2019年4月浙江选考第 23 题 10 分
%% typeId: 6
%% type: 计算题
%% chapter: 磁场
%% point: 带电粒子在组合场中的运动
%% method: 无
%% score: 10
%% degree: 350
%% duplicateId: 0
%% body:
有一种质谱仪由静电分析器和磁分析器组成，其简化原理如图所示。左侧静电分析器中有方向指向圆心 O、与 O 点等距离各点的场强大小相同的径向电场，右侧的磁分析器中分布着方向垂直于纸面向外的匀强磁场，其左边界与静电分析器的右边界平行，两者间距近似为零。离子源发出两种速度均为 $v_{0}$、电荷量均为 $q$、质量分别为 $m$ 和 $0.5m$ 的正离子束，从 M 点垂直该点电场方向进入静电分析器。在静电分析器中，质量为 $m$ 的离子沿半径为 $r_{0}$ 的四分之一圆弧轨道做匀速圆周运动，从 N 点水平射出，而质量为 $0.5m$ 的离子恰好从 ON 连线的中点 P 与水平方向成 $\theta$ 角射出，从静电分析器射出的这两束离子垂直磁场方向射入磁分析器中，最后打在放置于磁分析器左边界的探测板上，其中质量为 $m$ 的离子打在 O 点正下方的 Q 点。已知 $OP=0.5r_{0}$，$OQ=r_{0}$，N、P 两点间的电势差 $U_{NP}=\frac{mv^{2}}{q}$，$\cos\theta=\sqrt{\frac{4}{5}}$，不计重力和离子间相互作用。
\begin{enumerate}
	\item
	求静电分析器中半径为 $r_{0}$ 处的电场强度 $E_{0}$ 和磁分析器中的磁感应强度 $B$ 的大小；
	\item
	求质量为 $0.5m$ 的离子到达探测板上的位置与 O 点的距离 $l$（用 $r_{0}$ 表示）；
	\item
	若磁感应强度在（$B-\Delta B$）到（$B+\Delta B$）之间波动，要在探测板上完全分辨出质量为 $m$ 和 $0.5m$ 的两束离子，求 $\frac{{\Delta B}}{B}$ 的最大值。
\end{enumerate}
\onepicture[h=4.2cm, align=right]{figs/23.png}
%% answer: 
\jdanswer{
\begin{enumerate}
	\item
	$E_{0} = \frac{{mv_0^2}}{{q{r_0}}}$，$B = \frac{{m{v_0}}}{{q{r_0}}}$

	\item
	$1.5r_{0}$

	\item
	$12\%$
\end{enumerate}
}
%% memo:
\memoanswer{
（1）正离子束在电场中做圆周运动，由径向电场力提供向心力得 $qE_{0}=\frac{mv_{0}^{2}}{r_{0}}$，解得 $E_{0}=\frac{mv_{0}^{2}}{qr_{0}}$；进入磁场后，由洛伦兹力提供向心力得 $qv_{0}B=\frac{mv_{0}^{2}}{r_{0}}$，解得 $B=\frac{mv_{0}}{qr_{0}}$。

（2）质量为 $0.5m$ 的正离子束从开始运动到射出电场，由动能定理得 $qU_{NP}=\frac{1}{2}\times0.5mv^{2}-\frac{1}{2}\times0.5mv_{0}^{2}$，解得 $v=\sqrt{v_{0}^{2}+\frac{4qU_{NP}}{m}}=\sqrt{5}v_{0}$，故离子的运动半径为 $r=\frac{0.5mv}{qB}=\frac{1}{2}\sqrt{5}r_{0}$，由几何关系得 $l=2r\cos\theta-0.5r_{0}$，解得 $l=1.5r_{0}$。

（3）磁感应强度波动时，两束离子打到探测板上的位置不重合，临界情况：磁感应强度达到最大时，质量为 $0.5m$ 的离子打到的最远点，与磁感应强度为最小时，质量为 $m$ 的离子打到的最近点重合，由几何关系知，恰好能分辨的条件为 $2\cdot\frac{mv}{q(B-\Delta B)}-2\cdot\frac{0.5mv}{q(B+\Delta B)}=\frac{r_{0}}{2}$，即 $\frac{2r_{0}}{1-\frac{\Delta B}{B}}-\frac{2r\cos\theta}{1+\frac{\Delta B}{B}}=\frac{r_{0}}{2}$，解得 $\frac{\Delta B}{B}=\sqrt{17}-4=12\%$。
}

\end{enumerate}

\end{document}
